% GENERATED FILE. Edit canonical lessons, then run npm run generate:books.
% canonical-source-hash: fnv1a64:6a767367b5599193
% canonical-lessons: KA-C492-ullasa, KA-C492-anandisu, KA-C492-hummassu, KA-C492-harsa, KA-C492-trpti

\chapter{Joy and Delight}
\label{ch:ka-joy-and-delight}
\hlchaptermodality{492}

\begin{chapteropening}
\textbf{By the end of this chapter:} \emph{I can say cheerfulness, to enjoy, zest, delight and satisfaction.}

Complete the last lesson of chapter 492: I can say cheerfulness, to enjoy, zest, delight and satisfaction.
\end{chapteropening}

\section[\texorpdfstring{\kn{ಉಲ್ಲಾಸ} (ullāsa)}{ಉಲ್ಲಾಸ (ullāsa)}]{\kn{ಉಲ್ಲಾಸ} (ullāsa) — cheerfulness; high spirits}

\label{lesson:KA-C492-ullasa}



\begin{quote}
Before the new one: say the Kannada for regret, then the Kannada for remorse.
\end{quote}

\subsection*{You'll want to know: \kn{ಉಲ್ಲಾಸ}}
\textbf{\kn{ಉಲ್ಲಾಸ}} — \emph{ullāsa} — \textquotedblleft{}cheerfulness; high spirits\textquotedblright{}.

\begin{quote}
\kn{ಬೆಳಗಿನ} \kn{ತಣ್ಣನೆಯ} \kn{ಗಾಳಿಯಿಂದ} \kn{ಮನಸ್ಸಿಗೆ} \kn{ಉಲ್ಲಾಸ} \kn{ಆಯಿತು}.
\end{quote}

\begin{quote}
\emph{beḷagina taṇṇaneya gāḷiyinda manassige ullāsa āyitu.}
\end{quote}

\begin{quote}
The cool morning breeze lifted my spirits.
\end{quote}

\kn{ಉಲ್ಲಾಸದಿಂದ}, cheerfully.

\begin{cousinweb}
\textbf{\kn{ಉಲ್ಲಾಸ}} (\emph{ullāsa}) is borrowed from Sanskrit \emph{ullāsa}, \textquotedblleft{}light, splendour; a coming forth; joy, merriness\textquotedblright{} (Monier-Williams).
\end{cousinweb}

\begin{grammarlens}[title={What you've built}]
One more word for describing what you lived and felt.
\end{grammarlens}

\subsection*{Your turn}
\begin{itemize}
\raggedright
  \item \emph{Say it:} \emph{ullāsa}
  \item \emph{Say it:} \emph{ullāsa}, once more
  \item \emph{Say it:} \emph{beḷagina taṇṇaneya gāḷiyinda manassige ullāsa āyitu.} — The cool morning breeze lifted my spirits.
  \item \emph{Recall:} say the Kannada for regret, then the Kannada for remorse, then say \emph{ullāsa} again
\end{itemize}

\begin{tcolorbox}[breakable,colback=teal!4,colframe=teal!35!black,title={Before you move on}]
What does \kn{ಉಲ್ಲಾಸ} mean? (\textquotedblleft{}Cheerfulness; high spirits\textquotedblright{}.) Say it once more.
\end{tcolorbox}
\section[\texorpdfstring{\kn{ಆನಂದಿಸು} (ānandisu)}{ಆನಂದಿಸು (ānandisu)}]{\kn{ಆನಂದಿಸು} (ānandisu) — to enjoy}

\label{lesson:KA-C492-anandisu}



\begin{quote}
Before the new one: say the Kannada for remorse, then the Kannada for cheerfulness.
\end{quote}

\subsection*{You'll want to know: \kn{ಆನಂದಿಸು}}
\textbf{\kn{ಆನಂದಿಸು}} — \emph{ānandisu} — \textquotedblleft{}to enjoy\textquotedblright{}.

\begin{quote}
\kn{ಮಕ್ಕಳು} \kn{ಕಡಲತೀರದಲ್ಲಿ} \kn{ಆಟವನ್ನು} \kn{ಆನಂದಿಸಿದರು}.
\end{quote}

\begin{quote}
\emph{makkaḷu kaḍalatīradalli āṭavannu ānandisidaru.}
\end{quote}

\begin{quote}
The children enjoyed playing on the beach.
\end{quote}

Past (I …): \kn{ಆನಂದಿಸಿದೆ}, \emph{ānandiside}.

What you enjoy takes -\kn{ಅನ್ನು}. The ending -\kn{ಇಸು} makes a verb of the Sanskrit noun \kn{ಆನಂದ}, bliss.

\begin{cousinweb}
\textbf{\kn{ಆನಂದಿಸು}} (\emph{ānandisu}) is built from Sanskrit \emph{ānanda}, \textquotedblleft{}happiness\textquotedblright{} (CDIAL 1172), with the Kannada ending -\emph{isu} that turns a Sanskrit noun into a verb.
\end{cousinweb}

\begin{grammarlens}[title={What you've built}]
One more word for describing what you lived and felt.
\end{grammarlens}

\subsection*{Your turn}
\begin{itemize}
\raggedright
  \item \emph{Say it:} \emph{ānandisu}
  \item \emph{Say it:} \emph{ānandisu}, once more
  \item \emph{Say it:} \emph{makkaḷu kaḍalatīradalli āṭavannu ānandisidaru.} — The children enjoyed playing on the beach.
  \item \emph{Recall:} say the Kannada for remorse, then the Kannada for cheerfulness, then say \emph{ānandisu} again
\end{itemize}

\begin{tcolorbox}[breakable,colback=teal!4,colframe=teal!35!black,title={Before you move on}]
What does \kn{ಆನಂದಿಸು} mean? (\textquotedblleft{}To enjoy\textquotedblright{}.) Say it once more.
\end{tcolorbox}
\section[\texorpdfstring{\kn{ಹುಮ್ಮಸ್ಸು} (hummassu)}{ಹುಮ್ಮಸ್ಸು (hummassu)}]{\kn{ಹುಮ್ಮಸ್ಸು} (hummassu) — zest; eagerness}

\label{lesson:KA-C492-hummassu}



\begin{quote}
Before the new one: say the Kannada for cheerfulness, then the Kannada for to enjoy.
\end{quote}

\subsection*{You'll want to know: \kn{ಹುಮ್ಮಸ್ಸು}}
\textbf{\kn{ಹುಮ್ಮಸ್ಸು}} — \emph{hummassu} — \textquotedblleft{}zest; eagerness\textquotedblright{}.

\begin{quote}
\kn{ಹೊಸ} \kn{ಕೆಲಸ} \kn{ಶುರು} \kn{ಮಾಡಲು} \kn{ಅವನಿಗೆ} \kn{ತುಂಬಾ} \kn{ಹುಮ್ಮಸ್ಸು} \kn{ಇದೆ}.
\end{quote}

\begin{quote}
\emph{hosa kelasa śuru māḍalu avanige tumbā hummassu ide.}
\end{quote}

\begin{quote}
He is full of zest to start the new job.
\end{quote}

Close to \kn{ಉತ್ಸಾಹ}, but warmer and more everyday.

\begin{cousinweb}
\textbf{\kn{ಹುಮ್ಮಸ್ಸು}} (\emph{hummassu}) is an everyday Kannada word for eager energy; its origin is not settled.
\end{cousinweb}

\begin{grammarlens}[title={What you've built}]
One more word for describing what you lived and felt.
\end{grammarlens}

\subsection*{Your turn}
\begin{itemize}
\raggedright
  \item \emph{Say it:} \emph{hummassu}
  \item \emph{Say it:} \emph{hummassu}, once more
  \item \emph{Say it:} \emph{hosa kelasa śuru māḍalu avanige tumbā hummassu ide.} — He is full of zest to start the new job.
  \item \emph{Recall:} say the Kannada for cheerfulness, then the Kannada for to enjoy, then say \emph{hummassu} again
\end{itemize}

\begin{tcolorbox}[breakable,colback=teal!4,colframe=teal!35!black,title={Before you move on}]
What does \kn{ಹುಮ್ಮಸ್ಸು} mean? (\textquotedblleft{}Zest; eagerness\textquotedblright{}.) Say it once more.
\end{tcolorbox}
\section[\texorpdfstring{\kn{ಹರ್ಷ} (harṣa)}{ಹರ್ಷ (harṣa)}]{\kn{ಹರ್ಷ} (harṣa) — delight; joy}

\label{lesson:KA-C492-harsa}



\begin{quote}
Before the new one: say the Kannada for to enjoy, then the Kannada for zest.
\end{quote}

\subsection*{You'll want to know: \kn{ಹರ್ಷ}}
\textbf{\kn{ಹರ್ಷ}} — \emph{harṣa} — \textquotedblleft{}delight; joy\textquotedblright{}.

\begin{quote}
\kn{ಮದುವೆಯ} \kn{ಸುದ್ದಿಯಿಂದ} \kn{ಮನೆಯಲ್ಲಿ} \kn{ಹರ್ಷ} \kn{ತುಂಬಿತು}.
\end{quote}

\begin{quote}
\emph{maduveya suddiyinda maneyalli harṣa tumbitu.}
\end{quote}

\begin{quote}
The wedding news filled the house with joy.
\end{quote}

Written and formal; in talk people say \kn{ಸಂತೋಷ}.

\begin{cousinweb}
\textbf{\kn{ಹರ್ಷ}} (\emph{harṣa}) is borrowed from Sanskrit \emph{harṣa}, \textquotedblleft{}bristling, a thrill\textquotedblright{} (CDIAL 13999): the hair standing on end with excitement, and so delight.
\end{cousinweb}

\begin{grammarlens}[title={What you've built}]
One more word for describing what you lived and felt.
\end{grammarlens}

\subsection*{Your turn}
\begin{itemize}
\raggedright
  \item \emph{Say it:} \emph{harṣa}
  \item \emph{Say it:} \emph{harṣa}, once more
  \item \emph{Say it:} \emph{maduveya suddiyinda maneyalli harṣa tumbitu.} — The wedding news filled the house with joy.
  \item \emph{Recall:} say the Kannada for to enjoy, then the Kannada for zest, then say \emph{harṣa} again
\end{itemize}

\begin{tcolorbox}[breakable,colback=teal!4,colframe=teal!35!black,title={Before you move on}]
What does \kn{ಹರ್ಷ} mean? (\textquotedblleft{}Delight; joy\textquotedblright{}.) Say it once more.
\end{tcolorbox}
\section[\texorpdfstring{\kn{ತೃಪ್ತಿ} (tṛpti)}{ತೃಪ್ತಿ (tṛpti)}]{\kn{ತೃಪ್ತಿ} (tṛpti) — satisfaction; contentment}

\label{lesson:KA-C492-trpti}



\begin{quote}
Before the new one: say the Kannada for zest, then the Kannada for delight.
\end{quote}

\subsection*{You'll want to know: \kn{ತೃಪ್ತಿ}}
\textbf{\kn{ತೃಪ್ತಿ}} — \emph{tṛpti} — \textquotedblleft{}satisfaction; contentment\textquotedblright{}.

\begin{quote}
\kn{ಒಳ್ಳೆಯ} \kn{ಊಟದ} \kn{ನಂತರ} \kn{ಅಜ್ಜನಿಗೆ} \kn{ತೃಪ್ತಿ} \kn{ಆಯಿತು}.
\end{quote}

\begin{quote}
\emph{oḷḷeya ūṭada nantara ajjanige tṛpti āyitu.}
\end{quote}

\begin{quote}
After a good meal, grandfather felt satisfied.
\end{quote}

\kn{ತೃಪ್ತಿ} \kn{ಆಗು}, to feel satisfied; its opposite is \kn{ಅತೃಪ್ತಿ}.

\begin{cousinweb}
\textbf{\kn{ತೃಪ್ತಿ}} (\emph{tṛpti}) is borrowed whole from Sanskrit \emph{tṛpti}, \textquotedblleft{}satisfaction, contentment\textquotedblright{} (Monier-Williams).
\end{cousinweb}

\begin{grammarlens}[title={What you've built}]
One more word for describing what you lived and felt.
\end{grammarlens}

\subsection*{Your turn}
\begin{itemize}
\raggedright
  \item \emph{Say it:} \emph{tṛpti}
  \item \emph{Say it:} \emph{tṛpti}, once more
  \item \emph{Say it:} \emph{oḷḷeya ūṭada nantara ajjanige tṛpti āyitu.} — After a good meal, grandfather felt satisfied.
  \item \emph{Recall:} say the Kannada for zest, then the Kannada for delight, then say \emph{tṛpti} again
\end{itemize}

\begin{tcolorbox}[breakable,colback=teal!4,colframe=teal!35!black,title={Before you move on}]
What does \kn{ತೃಪ್ತಿ} mean? (\textquotedblleft{}Satisfaction; contentment\textquotedblright{}.) Say it once more.
\end{tcolorbox}
